% Options for packages loaded elsewhere
% Options for packages loaded elsewhere
\PassOptionsToPackage{unicode}{hyperref}
\PassOptionsToPackage{hyphens}{url}
\PassOptionsToPackage{dvipsnames,svgnames,x11names}{xcolor}
%
\documentclass[
  letterpaper,
  DIV=11,
  numbers=noendperiod]{scrartcl}
\usepackage{xcolor}
\usepackage{amsmath,amssymb}
\setcounter{secnumdepth}{5}
\usepackage{iftex}
\ifPDFTeX
  \usepackage[T1]{fontenc}
  \usepackage[utf8]{inputenc}
  \usepackage{textcomp} % provide euro and other symbols
\else % if luatex or xetex
  \usepackage{unicode-math} % this also loads fontspec
  \defaultfontfeatures{Scale=MatchLowercase}
  \defaultfontfeatures[\rmfamily]{Ligatures=TeX,Scale=1}
\fi
\usepackage{lmodern}
\ifPDFTeX\else
  % xetex/luatex font selection
\fi
% Use upquote if available, for straight quotes in verbatim environments
\IfFileExists{upquote.sty}{\usepackage{upquote}}{}
\IfFileExists{microtype.sty}{% use microtype if available
  \usepackage[]{microtype}
  \UseMicrotypeSet[protrusion]{basicmath} % disable protrusion for tt fonts
}{}
\makeatletter
\@ifundefined{KOMAClassName}{% if non-KOMA class
  \IfFileExists{parskip.sty}{%
    \usepackage{parskip}
  }{% else
    \setlength{\parindent}{0pt}
    \setlength{\parskip}{6pt plus 2pt minus 1pt}}
}{% if KOMA class
  \KOMAoptions{parskip=half}}
\makeatother
% Make \paragraph and \subparagraph free-standing
\makeatletter
\ifx\paragraph\undefined\else
  \let\oldparagraph\paragraph
  \renewcommand{\paragraph}{
    \@ifstar
      \xxxParagraphStar
      \xxxParagraphNoStar
  }
  \newcommand{\xxxParagraphStar}[1]{\oldparagraph*{#1}\mbox{}}
  \newcommand{\xxxParagraphNoStar}[1]{\oldparagraph{#1}\mbox{}}
\fi
\ifx\subparagraph\undefined\else
  \let\oldsubparagraph\subparagraph
  \renewcommand{\subparagraph}{
    \@ifstar
      \xxxSubParagraphStar
      \xxxSubParagraphNoStar
  }
  \newcommand{\xxxSubParagraphStar}[1]{\oldsubparagraph*{#1}\mbox{}}
  \newcommand{\xxxSubParagraphNoStar}[1]{\oldsubparagraph{#1}\mbox{}}
\fi
\makeatother

\usepackage{color}
\usepackage{fancyvrb}
\newcommand{\VerbBar}{|}
\newcommand{\VERB}{\Verb[commandchars=\\\{\}]}
\DefineVerbatimEnvironment{Highlighting}{Verbatim}{commandchars=\\\{\}}
% Add ',fontsize=\small' for more characters per line
\usepackage{framed}
\definecolor{shadecolor}{RGB}{241,243,245}
\newenvironment{Shaded}{\begin{snugshade}}{\end{snugshade}}
\newcommand{\AlertTok}[1]{\textcolor[rgb]{0.68,0.00,0.00}{#1}}
\newcommand{\AnnotationTok}[1]{\textcolor[rgb]{0.37,0.37,0.37}{#1}}
\newcommand{\AttributeTok}[1]{\textcolor[rgb]{0.40,0.45,0.13}{#1}}
\newcommand{\BaseNTok}[1]{\textcolor[rgb]{0.68,0.00,0.00}{#1}}
\newcommand{\BuiltInTok}[1]{\textcolor[rgb]{0.00,0.23,0.31}{#1}}
\newcommand{\CharTok}[1]{\textcolor[rgb]{0.13,0.47,0.30}{#1}}
\newcommand{\CommentTok}[1]{\textcolor[rgb]{0.37,0.37,0.37}{#1}}
\newcommand{\CommentVarTok}[1]{\textcolor[rgb]{0.37,0.37,0.37}{\textit{#1}}}
\newcommand{\ConstantTok}[1]{\textcolor[rgb]{0.56,0.35,0.01}{#1}}
\newcommand{\ControlFlowTok}[1]{\textcolor[rgb]{0.00,0.23,0.31}{\textbf{#1}}}
\newcommand{\DataTypeTok}[1]{\textcolor[rgb]{0.68,0.00,0.00}{#1}}
\newcommand{\DecValTok}[1]{\textcolor[rgb]{0.68,0.00,0.00}{#1}}
\newcommand{\DocumentationTok}[1]{\textcolor[rgb]{0.37,0.37,0.37}{\textit{#1}}}
\newcommand{\ErrorTok}[1]{\textcolor[rgb]{0.68,0.00,0.00}{#1}}
\newcommand{\ExtensionTok}[1]{\textcolor[rgb]{0.00,0.23,0.31}{#1}}
\newcommand{\FloatTok}[1]{\textcolor[rgb]{0.68,0.00,0.00}{#1}}
\newcommand{\FunctionTok}[1]{\textcolor[rgb]{0.28,0.35,0.67}{#1}}
\newcommand{\ImportTok}[1]{\textcolor[rgb]{0.00,0.46,0.62}{#1}}
\newcommand{\InformationTok}[1]{\textcolor[rgb]{0.37,0.37,0.37}{#1}}
\newcommand{\KeywordTok}[1]{\textcolor[rgb]{0.00,0.23,0.31}{\textbf{#1}}}
\newcommand{\NormalTok}[1]{\textcolor[rgb]{0.00,0.23,0.31}{#1}}
\newcommand{\OperatorTok}[1]{\textcolor[rgb]{0.37,0.37,0.37}{#1}}
\newcommand{\OtherTok}[1]{\textcolor[rgb]{0.00,0.23,0.31}{#1}}
\newcommand{\PreprocessorTok}[1]{\textcolor[rgb]{0.68,0.00,0.00}{#1}}
\newcommand{\RegionMarkerTok}[1]{\textcolor[rgb]{0.00,0.23,0.31}{#1}}
\newcommand{\SpecialCharTok}[1]{\textcolor[rgb]{0.37,0.37,0.37}{#1}}
\newcommand{\SpecialStringTok}[1]{\textcolor[rgb]{0.13,0.47,0.30}{#1}}
\newcommand{\StringTok}[1]{\textcolor[rgb]{0.13,0.47,0.30}{#1}}
\newcommand{\VariableTok}[1]{\textcolor[rgb]{0.07,0.07,0.07}{#1}}
\newcommand{\VerbatimStringTok}[1]{\textcolor[rgb]{0.13,0.47,0.30}{#1}}
\newcommand{\WarningTok}[1]{\textcolor[rgb]{0.37,0.37,0.37}{\textit{#1}}}

\usepackage{longtable,booktabs,array}
\newcounter{none} % for unnumbered tables
\usepackage{calc} % for calculating minipage widths
% Correct order of tables after \paragraph or \subparagraph
\usepackage{etoolbox}
\makeatletter
\patchcmd\longtable{\par}{\if@noskipsec\mbox{}\fi\par}{}{}
\makeatother
% Allow footnotes in longtable head/foot
\IfFileExists{footnotehyper.sty}{\usepackage{footnotehyper}}{\usepackage{footnote}}
\makesavenoteenv{longtable}
\usepackage{graphicx}
\makeatletter
\newsavebox\pandoc@box
\newcommand*\pandocbounded[1]{% scales image to fit in text height/width
  \sbox\pandoc@box{#1}%
  \Gscale@div\@tempa{\textheight}{\dimexpr\ht\pandoc@box+\dp\pandoc@box\relax}%
  \Gscale@div\@tempb{\linewidth}{\wd\pandoc@box}%
  \ifdim\@tempb\p@<\@tempa\p@\let\@tempa\@tempb\fi% select the smaller of both
  \ifdim\@tempa\p@<\p@\scalebox{\@tempa}{\usebox\pandoc@box}%
  \else\usebox{\pandoc@box}%
  \fi%
}
% Set default figure placement to htbp
\def\fps@figure{htbp}
\makeatother





\setlength{\emergencystretch}{3em} % prevent overfull lines

\providecommand{\tightlist}{%
  \setlength{\itemsep}{0pt}\setlength{\parskip}{0pt}}



 


\KOMAoption{captions}{tableheading}
\makeatletter
\@ifpackageloaded{caption}{}{\usepackage{caption}}
\AtBeginDocument{%
\ifdefined\contentsname
  \renewcommand*\contentsname{Table of contents}
\else
  \newcommand\contentsname{Table of contents}
\fi
\ifdefined\listfigurename
  \renewcommand*\listfigurename{List of Figures}
\else
  \newcommand\listfigurename{List of Figures}
\fi
\ifdefined\listtablename
  \renewcommand*\listtablename{List of Tables}
\else
  \newcommand\listtablename{List of Tables}
\fi
\ifdefined\figurename
  \renewcommand*\figurename{Figure}
\else
  \newcommand\figurename{Figure}
\fi
\ifdefined\tablename
  \renewcommand*\tablename{Table}
\else
  \newcommand\tablename{Table}
\fi
}
\@ifpackageloaded{float}{}{\usepackage{float}}
\floatstyle{ruled}
\@ifundefined{c@chapter}{\newfloat{codelisting}{h}{lop}}{\newfloat{codelisting}{h}{lop}[chapter]}
\floatname{codelisting}{Listing}
\newcommand*\listoflistings{\listof{codelisting}{List of Listings}}
\makeatother
\makeatletter
\makeatother
\makeatletter
\@ifpackageloaded{caption}{}{\usepackage{caption}}
\@ifpackageloaded{subcaption}{}{\usepackage{subcaption}}
\makeatother
\usepackage{bookmark}
\IfFileExists{xurl.sty}{\usepackage{xurl}}{} % add URL line breaks if available
\urlstyle{same}
\hypersetup{
  colorlinks=true,
  linkcolor={blue},
  filecolor={Maroon},
  citecolor={Blue},
  urlcolor={Blue},
  pdfcreator={LaTeX via pandoc}}


\author{}
\date{}
\begin{document}

\renewcommand*\contentsname{Table of contents}
{
\setcounter{tocdepth}{3}
\tableofcontents
}

\begin{Shaded}
\begin{Highlighting}[]
\FunctionTok{\# Common Problems and Solutions}

\NormalTok{Everyone gets errors in R. Errors are normal.}

\NormalTok{The important skill is learning how to read the message and try a fix.}

\FunctionTok{\#\# Problem 1: File Not Found}

\FunctionTok{\#\#\# What You Might See}

\InformationTok{\textasciigrave{}\textasciigrave{}\textasciigrave{}text}
\InformationTok{Error in file(file, "rt") : cannot open the connection}
\InformationTok{\textasciigrave{}\textasciigrave{}\textasciigrave{}}

\NormalTok{or:}

\InformationTok{\textasciigrave{}\textasciigrave{}\textasciigrave{}text}
\InformationTok{No such file or directory}
\InformationTok{\textasciigrave{}\textasciigrave{}\textasciigrave{}}

\FunctionTok{\#\#\# What It Usually Means}

\NormalTok{R cannot find your file.}

\FunctionTok{\#\#\# What To Check}

\NormalTok{Run:}


\NormalTok{::: \{.cell\}}

\InformationTok{\textasciigrave{}\textasciigrave{}\textasciigrave{}\{.r .cell{-}code\}}
\FunctionTok{getwd}\NormalTok{()}
\InformationTok{\textasciigrave{}\textasciigrave{}\textasciigrave{}}

\NormalTok{::: \{.cell{-}output .cell{-}output{-}stdout\}}

\InformationTok{\textasciigrave{}\textasciigrave{}\textasciigrave{}}
\InformationTok{[1] "/Users/ylim2/Documents/01\_Work/01\_MA213/R\_onboarding/r{-}onboarding{-}guide"}
\InformationTok{\textasciigrave{}\textasciigrave{}\textasciigrave{}}


\NormalTok{:::}
\NormalTok{:::}


\NormalTok{Then check whether }\InformationTok{\textasciigrave{}mydata.csv\textasciigrave{}}\NormalTok{ is inside that folder.}

\NormalTok{Also run:}


\NormalTok{::: \{.cell\}}

\InformationTok{\textasciigrave{}\textasciigrave{}\textasciigrave{}\{.r .cell{-}code\}}
\FunctionTok{file.exists}\NormalTok{(}\StringTok{"mydata.csv"}\NormalTok{)}
\InformationTok{\textasciigrave{}\textasciigrave{}\textasciigrave{}}

\NormalTok{::: \{.cell{-}output .cell{-}output{-}stdout\}}

\InformationTok{\textasciigrave{}\textasciigrave{}\textasciigrave{}}
\InformationTok{[1] TRUE}
\InformationTok{\textasciigrave{}\textasciigrave{}\textasciigrave{}}


\NormalTok{:::}
\NormalTok{:::}


\NormalTok{If the result is }\InformationTok{\textasciigrave{}FALSE\textasciigrave{}}\NormalTok{, R cannot find the file.}

\FunctionTok{\#\#\# Possible Fixes}

\SpecialStringTok{{-} }\NormalTok{Move }\InformationTok{\textasciigrave{}mydata.csv\textasciigrave{}}\NormalTok{ into your project folder.}
\SpecialStringTok{{-} }\NormalTok{Check the spelling of the file name.}
\SpecialStringTok{{-} }\NormalTok{Make sure the file ends in }\InformationTok{\textasciigrave{}.csv\textasciigrave{}}\NormalTok{.}
\SpecialStringTok{{-} }\NormalTok{Make sure you are working in the correct folder.}

\NormalTok{::: \{.callout{-}warning\}}
\FunctionTok{\#\# Common Mistake}

\InformationTok{\textasciigrave{}mydata.csv\textasciigrave{}}\NormalTok{ and }\InformationTok{\textasciigrave{}mydata (1).csv\textasciigrave{}}\NormalTok{ are different file names.}
\NormalTok{:::}

\FunctionTok{\#\# Problem 2: Package Installation Issues}

\NormalTok{Sometimes your instructor may ask you to install a package.}

\NormalTok{Example:}

\InformationTok{\textasciigrave{}\textasciigrave{}\textasciigrave{}r}
\FunctionTok{install.packages}\NormalTok{(}\StringTok{"ggplot2"}\NormalTok{)}
\InformationTok{\textasciigrave{}\textasciigrave{}\textasciigrave{}}

\FunctionTok{\#\#\# What Can Go Wrong?}

\SpecialStringTok{{-} }\NormalTok{Internet connection problems.}
\SpecialStringTok{{-} }\NormalTok{Typing the package name incorrectly.}
\SpecialStringTok{{-} }\NormalTok{R asks you to restart.}
\SpecialStringTok{{-} }\NormalTok{You do not have permission to install software on the computer.}

\FunctionTok{\#\#\# Possible Fixes}

\SpecialStringTok{{-} }\NormalTok{Check your internet connection.}
\SpecialStringTok{{-} }\NormalTok{Make sure the package name is inside quotation marks.}
\SpecialStringTok{{-} }\NormalTok{Restart RStudio and try again.}
\SpecialStringTok{{-} }\NormalTok{Ask your instructor or campus IT for help if you are on a lab computer.}

\NormalTok{::: \{.callout{-}tip\}}
\FunctionTok{\#\# Tip}

\NormalTok{You only need to install a package once on your computer.}

\NormalTok{But you need to load it with }\InformationTok{\textasciigrave{}library()\textasciigrave{}}\NormalTok{ each time you start a new R session.}
\NormalTok{:::}

\FunctionTok{\#\# Problem 3: Typing Mistakes}

\NormalTok{R is very picky about spelling.}

\NormalTok{These are all different to R:}

\InformationTok{\textasciigrave{}\textasciigrave{}\textasciigrave{}r}
\NormalTok{quiz\_score}
\NormalTok{Quiz\_score}
\NormalTok{quizscore}
\NormalTok{quiz.score}
\InformationTok{\textasciigrave{}\textasciigrave{}\textasciigrave{}}

\FunctionTok{\#\#\# Common Typing Mistakes}

\PreprocessorTok{|}\NormalTok{ Mistake }\PreprocessorTok{|}\NormalTok{ Example }\PreprocessorTok{|}\NormalTok{ Fix }\PreprocessorTok{|}
\PreprocessorTok{|{-}{-}{-}|{-}{-}{-}|{-}{-}{-}|}
\PreprocessorTok{|}\NormalTok{ Missing quote }\PreprocessorTok{|} \InformationTok{\textasciigrave{}read.csv(mydata.csv)\textasciigrave{}} \PreprocessorTok{|} \InformationTok{\textasciigrave{}read.csv("mydata.csv")\textasciigrave{}} \PreprocessorTok{|}
\PreprocessorTok{|}\NormalTok{ Misspelled object }\PreprocessorTok{|} \InformationTok{\textasciigrave{}quz\_score\textasciigrave{}} \PreprocessorTok{|} \InformationTok{\textasciigrave{}quiz\_score\textasciigrave{}} \PreprocessorTok{|}
\PreprocessorTok{|}\NormalTok{ Missing parenthesis }\PreprocessorTok{|} \InformationTok{\textasciigrave{}mean(data$quiz\_score\textasciigrave{}} \PreprocessorTok{|} \InformationTok{\textasciigrave{}mean(data$quiz\_score)\textasciigrave{}} \PreprocessorTok{|}
\PreprocessorTok{|}\NormalTok{ Wrong capitalization }\PreprocessorTok{|} \InformationTok{\textasciigrave{}Read.csv()\textasciigrave{}} \PreprocessorTok{|} \InformationTok{\textasciigrave{}read.csv()\textasciigrave{}} \PreprocessorTok{|}

\FunctionTok{\#\# Problem 4: Working Directory Problems}

\NormalTok{If R is looking in the wrong folder, it may not find your files.}

\NormalTok{Check the working directory with:}


\NormalTok{::: \{.cell\}}

\InformationTok{\textasciigrave{}\textasciigrave{}\textasciigrave{}\{.r .cell{-}code\}}
\FunctionTok{getwd}\NormalTok{()}
\InformationTok{\textasciigrave{}\textasciigrave{}\textasciigrave{}}

\NormalTok{::: \{.cell{-}output .cell{-}output{-}stdout\}}

\InformationTok{\textasciigrave{}\textasciigrave{}\textasciigrave{}}
\InformationTok{[1] "/Users/ylim2/Documents/01\_Work/01\_MA213/R\_onboarding/r{-}onboarding{-}guide"}
\InformationTok{\textasciigrave{}\textasciigrave{}\textasciigrave{}}


\NormalTok{:::}
\NormalTok{:::}


\NormalTok{In RStudio, using an RStudio Project can help keep your working directory organized.}

\NormalTok{::: \{.callout{-}tip\}}
\FunctionTok{\#\# Tip}

\NormalTok{For each assignment or project, keep your script and data file in the same folder.}
\NormalTok{:::}

\FunctionTok{\#\# Problem 5: Rendering Errors}

\NormalTok{When rendering a Quarto document, R runs the code from top to bottom.}

\NormalTok{A document may fail to render if:}

\SpecialStringTok{{-} }\NormalTok{A file is missing.}
\SpecialStringTok{{-} }\NormalTok{A package is not loaded.}
\SpecialStringTok{{-} }\NormalTok{An object is used before it is created.}
\SpecialStringTok{{-} }\NormalTok{There is a typing mistake in a code chunk.}

\FunctionTok{\#\#\# Example}

\NormalTok{This can fail:}

\InformationTok{\textasciigrave{}\textasciigrave{}\textasciigrave{}r}
\FunctionTok{mean}\NormalTok{(data}\SpecialCharTok{$}\NormalTok{quiz\_score)}
\InformationTok{\textasciigrave{}\textasciigrave{}\textasciigrave{}}

\NormalTok{if }\InformationTok{\textasciigrave{}data\textasciigrave{}}\NormalTok{ was never created earlier in the document.}

\FunctionTok{\#\#\# Fix}

\NormalTok{Make sure this line appears first:}

\InformationTok{\textasciigrave{}\textasciigrave{}\textasciigrave{}r}
\NormalTok{data }\OtherTok{\textless{}{-}} \FunctionTok{read.csv}\NormalTok{(}\StringTok{"mydata.csv"}\NormalTok{)}
\InformationTok{\textasciigrave{}\textasciigrave{}\textasciigrave{}}

\FunctionTok{\#\# Problem 6: Object Not Found}

\FunctionTok{\#\#\# What You Might See}

\InformationTok{\textasciigrave{}\textasciigrave{}\textasciigrave{}text}
\InformationTok{Error: object \textquotesingle{}data\textquotesingle{} not found}
\InformationTok{\textasciigrave{}\textasciigrave{}\textasciigrave{}}

\FunctionTok{\#\#\# What It Means}

\NormalTok{R does not know what }\InformationTok{\textasciigrave{}data\textasciigrave{}}\NormalTok{ is.}

\FunctionTok{\#\#\# Fix}

\NormalTok{Create the object first:}

\InformationTok{\textasciigrave{}\textasciigrave{}\textasciigrave{}r}
\NormalTok{data }\OtherTok{\textless{}{-}} \FunctionTok{read.csv}\NormalTok{(}\StringTok{"mydata.csv"}\NormalTok{)}
\InformationTok{\textasciigrave{}\textasciigrave{}\textasciigrave{}}

\NormalTok{Then run your analysis.}

\FunctionTok{\#\# Problem 7: Forgetting Quotation Marks}

\NormalTok{This is incorrect:}

\InformationTok{\textasciigrave{}\textasciigrave{}\textasciigrave{}r}
\FunctionTok{read.csv}\NormalTok{(mydata.csv)}
\InformationTok{\textasciigrave{}\textasciigrave{}\textasciigrave{}}

\NormalTok{This is correct:}

\InformationTok{\textasciigrave{}\textasciigrave{}\textasciigrave{}r}
\FunctionTok{read.csv}\NormalTok{(}\StringTok{"mydata.csv"}\NormalTok{)}
\InformationTok{\textasciigrave{}\textasciigrave{}\textasciigrave{}}

\NormalTok{File names must be in quotation marks.}

\FunctionTok{\#\# Troubleshooting Checklist}

\NormalTok{When something goes wrong, ask:}

\SpecialStringTok{1. }\NormalTok{Did I spell everything correctly?}
\SpecialStringTok{2. }\NormalTok{Did I use quotation marks around file names?}
\SpecialStringTok{3. }\NormalTok{Did I run the earlier lines of code?}
\SpecialStringTok{4. }\NormalTok{Is my file in the right folder?}
\SpecialStringTok{5. }\NormalTok{Does }\InformationTok{\textasciigrave{}file.exists("mydata.csv")\textasciigrave{}}\NormalTok{ return }\InformationTok{\textasciigrave{}TRUE\textasciigrave{}}\NormalTok{?}
\SpecialStringTok{6. }\NormalTok{Did I close all parentheses?}
\SpecialStringTok{7. }\NormalTok{Did I restart RStudio and try again?}

\NormalTok{::: \{.callout{-}note\}}
\FunctionTok{\#\# Encouragement}

\NormalTok{Errors do not mean you are bad at R.}

\NormalTok{Errors are part of learning. Read the message, try one fix, and run the code again.}
\NormalTok{:::}

\FunctionTok{\#\# Practice: Diagnose the Problem}

\NormalTok{::: \{.callout{-}note\}}
\FunctionTok{\#\# Practice}

\NormalTok{For each example, decide what the problem is.}

\FunctionTok{\#\#\# Example 1}

\InformationTok{\textasciigrave{}\textasciigrave{}\textasciigrave{}r}
\FunctionTok{read.csv}\NormalTok{(mydata.csv)}
\InformationTok{\textasciigrave{}\textasciigrave{}\textasciigrave{}}

\NormalTok{Question: What is missing?}

\FunctionTok{\#\#\# Example 2}

\InformationTok{\textasciigrave{}\textasciigrave{}\textasciigrave{}r}
\FunctionTok{mean}\NormalTok{(data}\SpecialCharTok{$}\NormalTok{quiz\_score}
\StringTok{\textasciigrave{}\textasciigrave{}\textasciigrave{}}

\AttributeTok{Question: What is missing?}

\AttributeTok{\#\#\# Example 3}

\StringTok{\textasciigrave{}\textasciigrave{}\textasciigrave{}}\NormalTok{r}
\NormalTok{Data }\OtherTok{\textless{}{-}} \FunctionTok{read.csv}\NormalTok{(}\StringTok{"mydata.csv"}\NormalTok{)}
\FunctionTok{head}\NormalTok{(data)}
\StringTok{\textasciigrave{}\textasciigrave{}\textasciigrave{}}

\AttributeTok{Question: Why might this fail?}
\AttributeTok{:::}

\AttributeTok{\#\# Answers}

\AttributeTok{1. Example 1 is missing quotation marks around the file name.}
\AttributeTok{2. Example 2 is missing a closing parenthesis.}
\AttributeTok{3. Example 3 may fail because }\StringTok{\textasciigrave{}}\NormalTok{Data}\StringTok{\textasciigrave{}}\AttributeTok{ and }\StringTok{\textasciigrave{}}\NormalTok{data}\StringTok{\textasciigrave{}}\AttributeTok{ are different names in R.}

\AttributeTok{\#\# Summary}

\AttributeTok{In this section, you learned how to troubleshoot common problems, including:}

\AttributeTok{{-} File not found errors.}
\AttributeTok{{-} Package installation issues.}
\AttributeTok{{-} Typing mistakes.}
\AttributeTok{{-} Working directory problems.}
\AttributeTok{{-} Rendering errors.}
\AttributeTok{{-} Object not found errors.}
\end{Highlighting}
\end{Shaded}





\end{document}
